\section{Evaluation protocol}
\label{sec:protocol}

The protocol, its hypotheses and its analysis were committed before the runs they govern;
every later amendment is dated in the repository, and no amendment changed an analysis
method after its data existed. Three analysis scripts received input-reading fixes after their
data existed (Section~\ref{sec:limitations}; Appendix~\ref{app:repro}).

\paragraph{Initial-state suites.} Each task has disjoint suites of
\rf{proto.nstates} serialized initial states drawn from robosuite's standard reset
distribution with reserved seed blocks: a \emph{validation} suite, which selects every
checkpoint and configuration; a \emph{held-out test} suite, locked in code until a freeze
document was committed and then run once per selected checkpoint; and a fresh
\emph{robustness base} suite (seeds from \rf{rob.seedbase}) for the perturbation study. No
two suites share a state or a seed.

\paragraph{Hypotheses.} H1: every trained family beats a uniform random-action baseline.
H2 (primary): diffusion versus BC-MLP and versus BC-RNN. H2-c: diffusion versus
BC-MLP-chunk, required before any diffusion advantage over BC-MLP is attributed to the
generative model rather than to chunking. H3: observation history, BC-RNN versus BC-MLP (and,
in ablation A3, diffusion with one versus two observations). H4: robustness under two
implemented perturbations (Section~\ref{sec:robustness}). H5 belongs to the separately frozen
point-cloud study (Section~\ref{sec:perception}) and is never pooled with H1 to H4. Each task
has \rf{proto.ncomparisons} declared success differences.

\paragraph{Selection.} Each training run's checkpoint is the one with the most validation
successes; ties go to the lower mean first-success step and then to the earlier checkpoint.
The same rule chose each family's tuning trial. Because the validation suite chose among
\rf{train.K} checkpoints and \rf{train.trials} trials, validation success is biased upward;
it is never reported as a test result (Table~\ref{tab:selection} in the appendix compares it
with success on fresh and held-out starts).

\paragraph{Rollouts.} Every start runs once ($R = 1$) with policy seed \rf{proto.policyseed}.
For the diffusion family the initial noise of each decision is a deterministic function of
the start, the policy seed and the decision index, so a rerun reproduces every action. With
$R = 1$ the difficulty of a start and the diffusion noise drawn on it cannot be separated
(Section~\ref{sec:limitations}).

\paragraph{Statistics.} The primary metric is the success rate. Families are compared by a
hierarchical paired bootstrap~\citep{efron1979bootstrap,field2007clustered} with
$B$\,=\,\rf{proto.B} replicates and seed \rf{proto.bootseed}: each replicate resamples the
\rf{proto.nseeds} training seeds of each family with replacement, then resamples the
\rf{proto.nstates} starts once and shares them across families, so every difference is
paired by start. A difference is \emph{separated} when its \rf{proto.level}\,\% percentile
interval excludes zero, and \emph{not separated} otherwise; not separated is not evidence of
equality. Intervals are not adjusted for multiplicity, and the counts of declared comparisons
are reported with each analysis. Per-checkpoint rates carry Wilson
intervals~\citep{wilson1927}. The secondary metric is the censored episode length: the
first-success step, or the horizon \rf{proto.horizon} for a failure. Because a failure
contributes the full horizon, a difference in censored length mixes a difference in success
with a difference in speed; we decompose it in Section~\ref{sec:results}.

\paragraph{Sensitivity analysis (not preregistered).} With \rf{proto.nseeds} seeds per family
the seed component of each bootstrap interval rests on very few values. After the audit
(Section~\ref{sec:limitations}) we recomputed every declared difference with a seed-level
$t$ interval on the frozen per-seed rates: Welch's interval for separately trained families,
a paired $t$ interval for the same checkpoints under two conditions, and Welch's interval on
the per-seed drops for a difference of drops. It treats per-seed means as the units, ignores
pairing by start, and is undefined when the seeds do not vary. It was not preregistered, it
changes no frozen verdict, and it is reported beside the frozen result as a check of how much
each verdict depends on the method.

\paragraph{Failure taxonomy.} Every failed episode is assigned, by a rule fixed before the
test and computed from recorded states, to one of seven failure categories: engineering
error, gripper never closed, no grasp with the gripper near the object, no grasp after
drifting away, object dropped, holding the object at timeout, and (Can) object in the goal
bin without success.
